\section{The Gornahoor Project}

Based on some exchanges lately, we feel the need to make clear the purpose of the \textbf{Gornahoor Project}. There is one distinction, and one distinction only: that between Tradition and anti-Tradition. Elaborating that distinction in its past and current manifestations is a legitimate topic for discussion. However, proselytizing for one religious form --- be it Buddhism, Catholicism, Asatru, Brahminism, or whatever --- is not legitimate.

A corollary to this is the distinction between principled and unprincipled. We are interested in principles only: that which is true always and everywhere. Hence, any arguments based on empirical, historical, personal, or other contingent factors are necessarily inconclusive.

Keep in mind the components of the Gornahoor motto:

\paragraph{Veritatem Loquor}


\emph{Veritatem loquor} means ``to speak the truth'', not ``to speak one's mind''. Not every, and more likely not many if any at all, thought that passes through the mind has any relation to Truth. So rather than expressing every passing thought, it would be better instead to trace it back to its roots. You will discover that ``your'' thoughts are not really yours at all.

\paragraph{Scientiam Acquirere}


\emph{Scientiam acquirere} means ``to acquire knowledge''. Of course, we mean here metaphysical knowledge. This is direct intuitive knowledge, beyond religion, science, and history. Although empirical knowledge of the latter fields may be used to illustrate and make points, they must never be confused with knowledge itself.

\paragraph{Liber Esse}


\emph{Liber esse} means ``to be free''. Unlike Rousseau, who claimed man is ``born free'', we claim that man must make himself free. He must transcend the spontaneous, determinative, and contingent factors that condition his Will in order to become conscious, creative, and free.


\flrightit{Posted on 2010-08-25 by Cologero}

\begin{center}* * *\end{center}

\begin{footnotesize}
\begin{sffamily}
\texttt{mleow on 2010-08-25 at 03:40 said:}

I recently found ``the gornahoor project'', and i am so far impressed by the content of this page and its authors, especially welcome is the translations of some more rare Evola passages.

It is also good that it avoids the currents of ``political traditionalism'' and ``the new right'', which usually follows in the traditional milieu, and instead focus solely on the spiritual path.

My only ``critic'' would be that it would be nice to see more lengthier and indepth postings, quality before quantity, right?

\hfill

\texttt{Sati Shankar on 2010-08-26 at 01:46 said:}

Yes it must adhere to its objectives.


\hfill

\end{sffamily}
\end{footnotesize}

